\section{Methodology}
\label{sec:methodology}

\subsection{What is evaluated}

\colbin{} is evaluated as a serialization format in Go: the bytes a Go
program writes and reads, against the Go implementations of the other
formats. One of its uses is a server's response to a web client in place of
JSON, but decoding in a browser is a different implementation and is outside
this paper; nothing here should be read as a claim about it.

\subsection{Datasets}
\label{sec:datasets}

All data is public and none of it was designed for this evaluation. Nothing
downloaded is stored in the repository: each dataset has a script that
downloads or generates it and checks every file against a manifest of SHA-256
checksums before anything reads it.

\begin{description}
\item[TPC-H]~\cite{tpch} at scale factor 0.1, generated by \code{dbgen}
  2.17.3: all eight tables, each as a case, and two composite cases, the
  whole schema as one message (\emph{normalized}) and orders with their line
  items nested inside them (\emph{nested}), the only parent and child the
  schema has.
\item[TPC-DS]~\cite{tpcds} at scale factor 0.1: all 24 tables, one case
  each.
\item[Public BI Benchmark]~\cite{publicbi}: the first table of each of the 46
  workbooks, and from each a seeded uniform random sample whose size depends
  on the table's place in the list: 350 rows of the first table by name, 700
  of the second, and so on in steps of 350, the last at $2^{14}$; or the
  whole table where it is shorter.
\item[SchemaStore]~\cite{schemastore}: the 27 documents of Viotti and
  Kinderkhedia's corpus~\cite{viotti2022benchmark}, each through a Go struct.
\end{description}

\paragraph{How the subsets are chosen.} Each subset follows a rule fixed in
the repository (\code{datasets/README.md}), not a selection of tables or rows.
For TPC-H and TPC-DS the rule is everything at one scale factor; scale factor
0.1 is about 100\,MB of text per benchmark, which every format holds in memory
at once, and a larger scale would add rows to the bulk regime only. For Public
BI the rule is the first table of every workbook, so that no workbook weighs
more than another; in 38 of the 46 workbooks it is the only table or the
others share its schema. The sizes are staggered so that the tables of web
size cover the whole range a response has, from a page of a few hundred
records to the regime's bound, rather than all sitting at that bound, where
a columnar layout is at its best; which table gets which size is decided by
its name alone. The rows are a random sample rather than the head of the file because
these tables are often sorted by their leading columns, so the head repeats
values more than the table does, which would flatter every format that
exploits repetition, \colbin{} most of all. The sample depends only on
fixed seeds, the table's name and the row index, and keeps table order: the
dataset's script draws $2^{14}$ rows of each table, and the case keeps a
second, independent draw of those. For
SchemaStore the selection is that of the corpus's authors.

\paragraph{Go types.} Every format is driven from one Go struct per dataset
root, written as a Go programmer would write it and not as any format
prefers: the narrowest integer type that holds the column per the source's
schema, money and decimals as integer minor units, dates as integer days, a
float only where the source has a float, and a pointer for a nullable column.
The schemas of the other formats (\code{.proto}, FlatBuffers, Cap'n Proto and
Avro schemas, generated MessagePack code, Arrow builders) are generated from
those structs. For TPC-H, TPC-DS and Public BI the structs themselves are
generated from the sources' \code{CREATE TABLE} statements. For SchemaStore
they were written by hand, from the Protocol Buffers schemas the corpus
provides; this is a place where the author's judgement enters
(Section~\ref{sec:limitations}).

\subsection{Kinds of message}
\label{sec:kinds}

We take $2^{14}$ rows as the bound of what a server sends to a web client in
one response; the bound is a convention of this paper, not a measurement.
Every case falls in one of two regimes by its row count alone:
\emph{web}, at most $2^{14} = 16\,384$ rows, and \emph{bulk}, more than that.
The web regime has \result{cases.web} cases: the \result{cases.web.publicbi}
Public BI samples, the \result{cases.web.schemastore} SchemaStore documents,
\result{cases.web.tpcds} TPC-DS tables and \result{cases.web.tpch} TPC-H
tables. The bulk regime has \result{cases.bulk}: the other
\result{cases.bulk.tpcds} TPC-DS tables, and for TPC-H four tables and the
two composite cases. The web regime holds two kinds of message that a
format treats differently, and the results keep them apart: a \emph{web
table}, a slice of rows (TPC-H, TPC-DS and Public BI;
\result{cases.categories.web-tables} cases), and a \emph{single record}, one
object with no table in it (SchemaStore;
\result{cases.categories.web-documents}). With \emph{bulk}
(\result{cases.categories.bulk}), these are the three kinds of message the
results are reported by, beside the four datasets, and the three uses the
format was built for (Section~\ref{sec:introduction}). SchemaStore's
documents are configuration files, not database records; they stand in for
the single record as the public corpus of small typed objects the
evaluation could use.

\subsection{Baselines}

Table~\ref{tab:baselines} lists the formats, their Go implementations and the
versions measured. Two formats appear with two implementations, because the
implementation moves the time by more than the format does: Protocol Buffers
with the reflection-based \fmt{protobuf-go} and with the generated code of
\fmt{vtprotobuf}~\cite{vtprotobuf}, and MessagePack with generated code and
with reflection. \colbin{} appears as three variants: the default
(\fmt{colbin}), with \fmt{packed5} on, and self-describing with its schema
section embedded (\fmt{colbin-json}). Only \fmt{colbin}, the default, is the result of this
paper; the others isolate one option each.

\begin{table}[t]
\centering\small
\resulttable{baselines}
\caption{Formats, implementations and versions. The versions are those pinned
by the artifact's \code{go.mod}; \colbin{}'s is the one the run recorded.}
\label{tab:baselines}
\end{table}

Not every format carries every case, and a format that cannot carry a case is
reported as not applicable, with the reason, rather than forced or silently
dropped (Table~\ref{tab:coverage}). Arrow carries only roots that are tables.
\colbin{} itself cannot carry
\result{coverage.web.colbin.na} of the 27 SchemaStore documents: it
refuses a struct with no field it can encode, and in those documents a type
has none. Cap'n Proto does not carry two SchemaStore documents, in which a
field's name collides with a method its Go code generator emits. \fmt{gob} does not transmit a pointer to a zero value, so it does
not round-trip any table with a nullable column that holds a zero; its sizes
on those cases are reported, flagged as lossy, and it is not timed on them
(Section~\ref{sec:limitations}).

\begin{table}[t]
\centering\small
\resulttable{coverage}
\caption{Cases per format and regime: round-tripped exactly, encoded but not
round-tripped (lossy; sized, not timed), and not applicable.}
\label{tab:coverage}
\end{table}

\subsection{What is measured}

\paragraph{Correctness first.} Before a size or a time is recorded, each
format encodes the case, decodes it, and the result is compared with the
original value by deep equality; where a case has an access aggregate, each
format's result is compared with the one computed on the Go value. A format
that fails a check on a case has no time for it.

\paragraph{Size.} The encoded message in bytes, as written and under two
general-purpose compressors at level 6: gzip~\cite{rfc1952}, which every
browser accepts and every server can apply to a response on the fly, and
Zstandard~\cite{rfc8878}. The text and the figures quote gzip; the tables
give both. Compressed sizes matter
because the realistic alternative to a binary format on the web is not JSON
but compressed JSON.

\paragraph{Time.} Four operations per format and case:
\emph{encode} into a reused buffer; \emph{decode-fresh}, into a value
allocated for the call; \emph{decode-reuse}, into one value held across
calls, which lets a decoder reuse its allocations; and \emph{access}, the sum
of one numeric column computed from the encoded bytes, in place where the
format allows it (Arrow, FlatBuffers, Cap'n Proto) and by decoding first
where it does not. Access is defined for the cases that have such a column
(\result{time.web.access.cases} of the web cases and
\result{time.bulk.access.cases} of the bulk cases).
For code-generated formats, converting between the format's own types and the
Go struct is outside the timed region.

\paragraph{Protocol.} Each case runs in its own process, pinned to one
hardware thread of a physical core with \code{taskset} and
\code{GOMAXPROCS=1}; the garbage collector is at its default setting, and
the runtime releases free memory to the kernel lazily
(\code{GODEBUG=madvdontneed=0}), for the reason given below. Web
cases run two at a time on two cores, bulk cases one at a time. Within a
case, every format and operation is checked and warmed up, the core is kept
busy for 200\,ms so that its clock is up before the first sample, and then
samples are taken, each after an untimed run of its own operation, round-robin over all formats and operations, so that each
benchmark sees the same neighbours and the same drift. A benchmark is
\result{machine.count} samples of at least \result{machine.sample} each, and
its value is their median. Sizes do not depend on the machine and are
computed in a separate stage.

\paragraph{Memory handed back between samples.} A process holds every
format's values of a case at once, and between two samples Go's runtime
returns free heap to the kernel. With its default on Linux
(\code{MADV\_DONTNEED}) the next benchmark that allocates takes a page
fault for every page it gets back, so what a decode into a fresh value cost
depended on which benchmark ran before it, and the round-robin order, being
fixed, repeated that in every round and in both passes. A first version of
this evaluation was measured that way, and the same decoder on the same
bytes, \colbin{}'s default and its variants, differed by up to a half on a
bulk case in decode-fresh while agreeing in decode-reuse, which allocates
nothing. Two things remove it. With \code{MADV\_FREE} the kernel takes
the pages only when it is short of memory, and the faults do not occur.
And every sample is preceded by an untimed run of the same operation for the
same time, followed by a collection, so that the memory a sample allocates
from is memory its own format has just used and freed, not another's.
\colbin{}'s variants, which share a decoder,
are in the timing tables of Appendix~\ref{sec:against-json} as the control
of this.

\paragraph{Summaries.} The results compare \colbin{} with each other format
directly: on each case both carry, \colbin{}'s number is divided by the
other format's for the same case and operation. A summary over cases is the
geometric mean of those ratios, in which every case weighs the same, with
the number of cases \colbin{} wins outright beside it. It is printed as a
factor, the inverse of that mean where \colbin{} wins ($2.5\times$ smaller,
or faster, for a mean of 0.4) and the mean itself, marked \emph{worse},
where it loses; the mean of the inverses is the inverse of the mean, so
this is the same statistic read the other way. The cases are
grouped two ways, each case in exactly one group: by dataset, and by kind of
message (Section~\ref{sec:kinds}). Because the comparison is always on the
cases both formats carry, a format that carries few cases is compared with
\colbin{} on those cases only. Each format against JSON, over the cases the
format carries and by regime, is in Appendix~\ref{sec:against-json}; there,
for access, which JSON does not have, the reference is JSON's decode-fresh,
what a client without it does instead. A message of zero bytes (a document
whose values are all defaults is zero bytes in \colbin{} and Protocol
Buffers) has a ratio of zero, which a geometric mean cannot hold; such cases
are left out of the mean and marked.

\FloatBarrier
\subsection{Machine and repeatability}
\label{sec:machine}

{\sloppy The main numbers were measured in Oracle Cloud Infrastructure container
instances of shape \code{CI.Standard.E4.Flex}, 4 OCPUs (physical cores
dedicated to the instance) and 32\,GB, on an \result{machine.cpu},
\result{machine.kernel}, Go \result{machine.go-version}, on CPUs
\result{machine.measured-cores}. Each dataset was measured in its own
instance, all four of the same shape and at the same commit, side by side;
every case, and so every ratio between two formats, comes from a single
process on a single core. Each instance was started from a pinned container
image, cloned this artifact at the commit it reports, and fetched and
verified its dataset itself.\par}

A container cannot fix the clock: the frequency governor and the boost are
the host's, and how far a core boosts depends on what the host's other
tenants do. We therefore measured the clock instead of fixing it. Before
every round of a case's samples, the harness times a chain of dependent
64-bit multiplies on the measured core, which runs at one multiply per three
cycles on this processor. Over the first web pass the clock was between
\result{probe.web.mhz.p5} and \result{probe.web.mhz.p95}\,MHz (5th and 95th
percentiles; median \result{probe.web.mhz.median}). The container exposes no other measure of the
clock: no frequency driver, no APERF/MPERF counters, a time-stamp counter
that runs at the nominal frequency, and no access to the cycle counter. The
absolute clock rests on the multiplier's documented latency; the variation,
which is what the ratios depend on, does not. Within a case, which is
what a ratio between two formats depends on, the slowest probe was within
\result{probe.web.spread.median}\,\% of the fastest for the median case and
within \result{probe.web.spread.p90}\,\% for nine cases in ten;
\result{probe.web.spread.over5} of \result{probe.web.cases} cases saw the
clock move by more than 5\,\%. Because samples are taken round-robin, such a
change falls on every format of the case alike. The probes, the machine
description and the per-second log of the cores are part of the artifact.

\paragraph{A second machine.} The same commit was measured again, in full, on
a machine of another kind: a laptop whose processor is a \result{laptop.machine.cpu}, a
15\,W client processor with two performance cores and eight efficiency
cores, \result{laptop.machine.kernel}, Go \result{laptop.machine.go-version}.
There the clock could be fixed rather than measured. Boost was disabled and
the \code{performance} governor set, so that the performance cores ran at
their base frequency throughout; the cases ran on one hardware thread of each
performance core (CPUs \result{laptop.machine.measured-cores}), with the
same protocol as above, and every other process and every interrupt were
confined to the efficiency cores. On this machine the processor's cycle
counter can be read, and it confirms the probe: over the first web pass the
probe read a median of \result{laptop.probe.web.mhz.median}\,MHz and the
cycle counter \result{laptop.probe.web.counted-mhz.median}\,MHz, at
\result{laptop.probe.web.cycles-per-multiply} cycles per multiply. Within a
case the probe moved by at most \result{laptop.probe.web.spread.p90}\,\% for
nine cases in ten, and by more than 5\,\% in
\result{laptop.probe.web.spread.over5} of \result{laptop.probe.web.cases}.
The two machines are reported apart and never pooled
(Section~\ref{sec:second-machine}); a number that does not say otherwise is
the cloud's.

The web cases were timed twice, several hours apart, to measure how far a
number moves when nothing changes. Of \result{repeat.lines} benchmarks,
\result{repeat.within3}\,\% differ by at most 3\,\% between the two passes,
\result{repeat.within5}\,\% by at most 5\,\% and \result{repeat.within10}\,\%
by at most 10\,\%; the median difference is \result{repeat.median}\,\%, the
90th percentile \result{repeat.p90}\,\%, and neither pass is slower than
the other overall. Absolute times on
this machine are therefore not repeatable to 3\,\%. Ratios between formats
are: the geometric means of each format's
time over \colbin{}'s differ between the two passes by at most
\result{repeat.ratio-worst}\,\% (Table~\ref{tab:repeatability}), because the
round-robin sampling exposes every format in a case to the same conditions.
For that reason this paper reports ratios, not absolute times, and treats a
difference between two formats smaller than about 4\,\% as no difference.
Timing tables use the first pass. On the laptop, with its clock fixed,
\result{laptop.repeat.within3}\,\% of \result{laptop.repeat.lines} benchmarks
differ by at most 3\,\% between the two passes, the median difference is
\result{laptop.repeat.median}\,\%, and the geometric means over \colbin{}'s
differ by at most \result{laptop.repeat.ratio-worst}\,\%.

\begin{table}[t]
\centering\small
\resulttable{repeatability}
\caption{Each format's time over \colbin{}'s, as a geometric mean over the
web cases both have, in the two timing passes.}
\label{tab:repeatability}
\end{table}
